\begin{center}
    \Large{\textbf{Аннотация}}
\end{center}

\begin{singlespace}

Решение уравнения частных производных находит применение в задачах математической физики.
В работе рассматривается их применение для решении задачи ЭЭГ. В задаче ЭЭГ трудность 
вызывает получение точного сигнала от головного мозга. Исследователи встречаются со следующими проблемами. 
Высокая чувствительность прибора к движениям и тремору, обусловленному психоэмоциональным напряжением пациента, 
вызывает помехи в работе, что может затруднить диагностику. 
А если встраивать датчики напрямую в кору мозга, то они уже сами начинают влиять на его работу, 
а через некоторое время приходят в негодность из-за перестройки мозга вокруг них. Поэтому 
стоит задача точного определения восстановления местоположения источников сигнала, 
при условии высокого уровня шума для работы с этими источниками в медицинских целях.
Целью работы является предложить метод решения обратной задачи в частных производных.
Предлагается использовать физико-информированный подход в восстановлении источников, 
использующийся в задачах восстановления временных рядов, вносящий априорные знания 
о модели для уменьшения уровня шума от данных.
\end{singlespace}
